% ---------------------------------------------------------------------
% Argomento: rete mista (R1 serie con R2//R3 serie con R4); si conosce la
%            corrente in R2 e si risale alla tensione E del generatore
% Livello: intermedio   |   Tempo stimato: 18 min
% Lezioni: 7/10 (serie e parallelo), 9/10 (KCL e KVL)
% ---------------------------------------------------------------------
\begin{esercizio}
\dato{R1}{14}{4}{5}               % [ohm]
\dato{R2}{91}{66}{49}
\dato{R3}{61}{99}{84}
\dato{R4}{3}{5}{11}
\dato{I2}{1.7}{2.8}{3.2}          % corrente in R2 [A]
\calcola[V_{23}]{V23}{\x{R2}*\x{I2}}
\calcola[I_3]{I3}{\x{V23}/\x{R3}}
\calcola[I]{I}{\x{I2}+\x{I3}}
\calcola[V_1]{V1}{\x{R1}*\x{I}}
\calcola[V_4]{V4}{\x{R4}*\x{I}}
\calcola[E]{E}{\x{V1}+\x{V23}+\x{V4}}

Sono note le resistenze e la corrente $I_2$ in $R_2$; $E$ non è nota.

\begin{center}
\begin{minipage}{0.52\linewidth}\centering
  \begin{circuitikz}[x=0.85cm, y=0.5cm]
    \ctikzset{bipoles/length=1cm}
    \draw (0,0) to[vsource, l=$E$, n=G] (0,3)
          to[R, l=$R_1$] (3,3) node[circ]{} -- (5.5,3)
          (3,3) to[R, l_=$R_2$] (3,0) node[circ]{}
          (5.5,3) to[R, l=$R_3$] (5.5,0) -- (3,0)
          (3,0) to[R, l=$R_4$] (0,0);
    \draw[-{Stealth[length=2mm]}] (3.5,2.4) -- (3.5,0.6) node[midway, right]{$I_2$};
    \morsetti{G}
  \end{circuitikz}
\end{minipage}\hfill
\begin{minipage}{0.4\linewidth}
  \renewcommand{\arraystretch}{1.1}
  \begin{tabular}{r@{${}={}$}l}
    $R_1$ & \mis{R1}{\ohm}\\
    $R_2$ & \mis{R2}{\ohm}\\
    $R_3$ & \mis{R3}{\ohm}\\
    $R_4$ & \mis{R4}{\ohm}\\
    $I_2$ & \mis{I2}{\ampere}
  \end{tabular}
\end{minipage}
\end{center}

\vspace{-6pt}
\begin{parts}
  \part[1] Calcola la tensione $V_{23}$ ai capi di $R_2$ (e quindi di $R_3$).
  \part[1] Calcola la corrente $I_3$ in $R_3$.
  \part[1] Calcola la corrente $I$ erogata dal generatore.
  \part[2] Calcola le tensioni $V_1$ e $V_4$ ai capi di $R_1$ e $R_4$.
  \part[1] Applica la seconda legge di Kirchhoff alla maglia $E$, $R_1$, $R_2$, $R_4$.
  \part[2] Calcola la tensione $E$ del generatore.
\end{parts}

\begin{solution}
  $V_{23} = R_2\,I_2 \approx \mis{V23}{\volt}$\\[4pt]
  $I_3 = V_{23}/R_3 \approx \mis{I3}{\ampere}$\\[4pt]
  KCL: $I = I_2 + I_3 \approx \mis{I}{\ampere}$\\[4pt]
  $V_1 = R_1\,I \approx \mis{V1}{\volt}$,\qquad $V_4 = R_4\,I \approx \mis{V4}{\volt}$\\[4pt]
  KVL sulla maglia $E$, $R_1$, $R_2$, $R_4$: $E = V_1 + V_{23} + V_4$\\[2pt]
  $E \approx \mis{E}{\volt}$
  \criterio{a) 1 pt (R2 e R3 hanno la stessa tensione); b) 1 pt; c) 1 pt KCL;
    d) 1 pt per ciascuna tensione; e) 1 pt equazione corretta $E=V_1+V_{23}+V_4$ (segni coerenti); f) 1 pt somma, 1 pt risultato con unità.
    Unità mancanti: $-0{,}5$ pt una volta sola.}
\end{solution}
\end{esercizio}
